\unit{Atomic spectroscopy, multipoles and radiative lifetimes}
\unitref{Week 9; Week 10; Assignment 4 Q3}
\subsection{Absorption and emission connect energy levels}
An atom absorbs light when an electromagnetic
mode drives a transition between its
quantum energy eigenstates. Energy
conservation fixes
\[
\hbar\omega_{if}=E_i-E_f,\quad
\nu=\frac{E_i-E_f}{h},\quad
\lambda=\frac{hc}{E_i-E_f}.
\]
A positive energy difference alone does
not guarantee a transition:
the light--matter interaction operator
must have a nonzero matrix element
between the two states.

\subsection{Electric dipole coupling}
The leading long-wavelength interaction
is \(H_{\rm int}=-\mathbf d\cdot\mathbf E\),
where the electron dipole is \(\mathbf d=-e\mathbf r\).
The transition amplitude contains
\[
\mathcal M_{fi}\propto
\bra f\mathbf d\cdot\boldsymbol\epsilon\ket i.
\]
This separates into radial and angular
factors when orbitals are expressed as
\(R_{nl}(r)Y_{lm}(\theta,\phi)\).
For an electric dipole the radial integral is
\[
\boxed{\int_0^\infty
R_f^*(r)R_i(r)\,r^3dr}
\]
or, using \(u=rR\),
\(\int u_f^*(r)\,r\,u_i(r)dr\).
The power \(r^3\) combines the
\(r^2dr\) volume measure with
the additional \(r\) from the dipole.

\subsection{Multipole rank and parity rules}
An electric multipole \(Ek\)
is a rank-\(k\) spherical tensor
of parity \((-1)^k\);
a magnetic multipole \(Mk\)
has parity \((-1)^{k+1}\).
For total angular momenta \(J,J'\),
the geometric rule is
\[
\boxed{|J-J'|\le k\le J+J',\qquad
m'=m+q,\quad |q|\le k.}
\]
For the most common transitions:
\[
\begin{array}{l|l|l}
&\text{Parity}&\text{Common }J\text{ rules}\\\hline
E1\ (k=1)&\text{changes}&
\Delta J=0,\pm1,\;0\nleftrightarrow0\\
M1\ (k=1)&\text{same}&
\Delta J=0,\pm1,\;0\nleftrightarrow0\\
E2\ (k=2)&\text{same}&
|\Delta J|\le2,\ J+J'\ge2
\end{array}
\]
The usual E1 orbital rule for one electron
is \(\Delta l=\pm1\).
For E2 the single-electron orbital rules
are \(\Delta l=0,\pm2\), subject to triangle
conditions. Thus \(s\leftrightarrow d\)
can be E2, but \(s\leftrightarrow p\)
is not E2 because parity changes.
Additional approximate LS coupling rules,
such as \(\Delta S=0\) for electric-dipole
transitions, may be softened by spin--orbit
mixing. A weak intercombination line can
therefore exist even if it is forbidden
in ideal LS coupling.

\begin{worked}{Which sodium transitions can radiate by E1?}
The valence-orbital transitions
\(3p\to3s\), \(4s\to3p\)
and \(3d\to3p\) all have
\(\Delta l=\pm1\), so their
orbital E1 matrix elements may
be nonzero. But \(3d\to3s\)
has \(\Delta l=-2\) and is E1 forbidden.
Always check \emph{downward energy ordering}
for spontaneous emission: an angularly
allowed upward transition is absorption,
not spontaneous decay from a lower level.
\end{worked}

\subsection{Why some states decay quickly}
An excited atom couples to vacuum
electromagnetic modes and can
spontaneously emit a photon. For an
E1 transition between appropriately
specified states,
\[
\boxed{A_{if}=
\frac{\omega_{if}^3}{3\pi\epsilon_0\hbar c^3}
|\bra f\mathbf d\ket i|^2,}
\]
where the displayed vector dipole
strength and sublevel/polarisation
summations must be handled consistently.
The strong \(\omega^3\) dependence means
larger energy gaps generally radiate
faster for comparable matrix elements.
A forbidden E1 line may instead decay
by much weaker M1/E2 channels,
making it useful for long-lived
clock or metastable states.

\subsection{Lifetime and branching ratios}
If an excited state can decay into
several lower states,
\[
\boxed{\Gamma_i=\sum_{f<i}A_{if},\quad
\tau_i=\Gamma_i^{-1},\quad
b_{if}=\frac{A_{if}}{\Gamma_i}.}
\]
Do \emph{not} use the inverse
of just one \(A\) if multiple
important decay channels exist.
In quantum optics \(\Gamma_i\)
is a population decay rate with units
\(\mathrm{s}^{-1}\); a natural
linewidth in angular-frequency units
is of order \(\Gamma_i\).

\begin{worked}{A transition-rate calculation with a stated dipole}
For a \emph{hypothetical} 589 nm line with
\(|\mathbf d|=2ea_0\) and no other
downward channels,
\[
\omega=\frac{2\pi c}{589\,\mathrm{nm}},
\quad
A=\frac{\omega^3(2ea_0)^2}
{3\pi\epsilon_0\hbar c^3}.
\]
This gives a rate of order
\(4\times10^7\ \mathrm{s}^{-1}\)
and a lifetime of order
\(25\) ns. It is an illustration
of using the formula, \emph{not}
a claim about a measured Na lifetime.
For actual sodium, compute the
wavefunction-derived matrix elements
and sum all allowed channels.
\end{worked}

\subsection{Einstein coefficients and induced processes}
An incident radiation field can
stimulate both absorption and emission.
Einstein's \(A\) coefficient describes
spontaneous decay; the \(B\) coefficients
describe rates proportional to the
radiation energy density.
Their detailed-balance relations explain
thermal equilibrium radiation and
underlie laser gain calculations.
A large population inversion makes
stimulated emission dominate absorption
at the selected lasing transition.

\begin{checkpoint}
A complete transition question requires
\emph{energy}, \emph{tensor rank},
\emph{angular-momentum triangle},
\emph{parity}, and a \emph{reduced
matrix element}. A lifetime then needs
all allowed downward branches.
\end{checkpoint}
